% Teaching manuscript. Edit this TeX source; Markdown is compatibility navigation only.

\StudyChapter{CS5222-Assignment02}{Assignment 2 · Chapters 3, 4 and 5｜rdt3.0、GBN 窗口推演与 OpenFlow 流表}
\StudyTarget{CS5222-Assignment02__root-section}

\StudyRepoRef{CS5222/course-notes/README.md}{}{课程目录} · \StudyRef{CS5222-Chapter03}{root}{Chapter 3} · \StudyRef{CS5222-Chapter04}{root}{Chapter 4} · \StudyRef{CS5222-Chapter05}{root}{Chapter 5} · \StudyRef{CS5222-Assignment01}{root}{上一份：Assignment 1} · \StudyRef{CS5222-Assignment03}{root}{下一份：Assignment 3}

\StudyHeading{1}{导读｜状态机画法、GBN 边界严密推演与 SDN 流表配置}
\StudyTarget{CS5222-Assignment02__overview}

Assignment 2（\texttt{CS5222 Assignment 2.pdf}）共五道经典教材习题：Chapter 3 部分的 \textbf{P8} 要求画出 \texttt{rdt3.0} 接收端有限状态机，\textbf{P22} 在无乱序信道上严格推导 GBN 发送窗口与在途确认号的全部可能取值（特别注意 \texttt{ACK(k-5)} 为何仍可在途）；Chapter 4 \& 5 部分的 \textbf{P4} 揭示传统纯目的地址转发表的局限，并自然引出 \textbf{P19–P20} 在三角形三交换机网络中编写 OpenFlow 广义流表实现双向互通与顺/逆时针按源分流。全部结论与参考标答逐项核对一致，配图均为原创绘制原理图。


\StudyHeading{1}{Ch.3 P8. 可靠传输协议 \texttt{rdt3.0} 接收端有限状态机（FSM）}
\StudyTarget{CS5222-Assignment02__p8}

\textbf{Original Question:} Draw the FSM for the receiver side of protocol \texttt{rdt3.0}.

\textbf{核心设计逻辑：}相比仅处理比特差错的 \texttt{rdt2.2}，\texttt{rdt3.0} 新增的挑战是底层信道可能丢包；但\textbf{丢包超时检测与重传职责完全落在发送端}（靠发送端的倒计时器完成）。在接收端看来，数据包在途中丢失与发送端暂时还没发新包没有任何区别——接收端只需老老实实停在当前状态等待期待序号（0 或 1）的无损分组，一旦收到期望分组就拆包上交并发回同号 ACK；一旦收到比特损坏分组或重复的上一号分组，就沿自环重发上一个已收妥分组的确认号（无 NAK 机制）。因此 \textbf{\texttt{rdt3.0} 接收端 FSM 与 \texttt{rdt2.2} 接收端 FSM 完全一致}，包含两个对称状态：

\begin{StudyCode}
Initial State -> S0: Wait for 0 from below
  [Transition S0 -> S1]
    Event : rdt_rcv(rcvpkt) && notcorrupt(rcvpkt) && has_seq0(rcvpkt)
    Action: extract(rcvpkt, data); deliver_data(data);
            sndpkt = make_pkt(ACK, 0, checksum); udt_send(sndpkt);

  [Self-loop S0 -> S0]
    Event : rdt_rcv(rcvpkt) && (corrupt(rcvpkt) || has_seq1(rcvpkt))
    Action: sndpkt = make_pkt(ACK, 1, checksum); udt_send(sndpkt);

State S1: Wait for 1 from below
  [Transition S1 -> S0]
    Event : rdt_rcv(rcvpkt) && notcorrupt(rcvpkt) && has_seq1(rcvpkt)
    Action: extract(rcvpkt, data); deliver_data(data);
            sndpkt = make_pkt(ACK, 1, checksum); udt_send(sndpkt);

  [Self-loop S1 -> S1]
    Event : rdt_rcv(rcvpkt) && (corrupt(rcvpkt) || has_seq0(rcvpkt))
    Action: sndpkt = make_pkt(ACK, 0, checksum); udt_send(sndpkt);
\end{StudyCode}

\textbf{English answer:} Because packet-loss recovery via timeout is handled entirely at the sender, the \texttt{rdt3.0} receiver FSM is identical to the \texttt{rdt2.2} receiver FSM: two states (\texttt{Wait for 0 from below} and \texttt{Wait for 1 from below}), advancing state and sending \texttt{ACK0}/\texttt{ACK1} on an uncorrupted expected packet, and self-looping while re-sending the previous \texttt{ACK1}/\texttt{ACK0} whenever a corrupted or duplicate packet arrives.


\StudyHeading{1}{Ch.3 P22. GBN 发送窗口集合与在途 ACK 取值范围（$N=4$，模 1024）}
\StudyTarget{CS5222-Assignment02__p22}

\textbf{Original Question:} Consider the GBN protocol with a sender window size of $4$ and a sequence number range of $1{,}024$. Suppose that at time $t$, the next in-order packet that the receiver is expecting has a sequence number of $k$. Assume that the medium does not reorder messages. Answer the following questions:
\begin{itemize}
\item \textbf{(a)} What are the possible sets of sequence numbers inside the sender's window at time $t$? Justify your answer.
\item \textbf{(b)} What are all possible values of the ACK field in all possible messages currently propagating back to the sender at time $t$? Justify your answer.
\end{itemize}

\textbf{严密上下界推导（所有序号均在 $\bmod\, 1024$ 意义下理解）：}
\begin{enumerate}
\item \textbf{(a) 时刻 $t$ 发送方窗口共有 5 种可能的连续 4 元集合}：
  既然接收方在时刻 $t$ 正期待按序分组 $k$，说明接收方\textbf{已经按序收妥了分组 $k-1$ 以及更早的所有分组}，并且已经依次向回程链路发出过 $\mathrm{ACK}(k-N), \dots, \mathrm{ACK}(k-1)$，但尚未收到分组 $k$。
  \begin{itemize}
  \item \textbf{最靠右情形}：若这批确认（或至少累积确认 $\mathrm{ACK}(k-1)$）已经在时刻 $t$ 之前抵达发送方，发送方窗口起点已推进至 $k$，窗口为 $[k,\; k+1,\; k+2,\; k+3]$。
  \item \textbf{最靠左情形}：既然接收方已经收到了分组 $k-1$，发送方在过去某时刻必定已经将 $k-1$ 发出；在窗口长度 $N = 4$ 的限制下，发送方必须早已收到 $\mathrm{ACK}(k-5)$（即 $k-N-1$），其窗口起点 $\mathrm{base}$ 至少推进到了 $(k-1) - (4-1) = k - 4$，才能把 $k-1$ 纳入发送窗口。若随后发回的 $\mathrm{ACK}(k-4), \dots, \mathrm{ACK}(k-1)$ 在时刻 $t$ 全都还在半路传播（或途中丢失），发送方窗口便仍停留在 $[k-4,\; k-3,\; k-2,\; k-1]$。
  \end{itemize}
  因此窗口起点可落在 $k-4$ 到 $k$ 之间的任意位置，共 \textbf{5 种可能窗口}：
  \[
  [k-4,\,k-1],\quad [k-3,\,k],\quad [k-2,\,k+1],\quad [k-1,\,k+2],\quad [k,\,k+3] \pmod{1024}.
  \tag{A2.1}
  \]

\item \textbf{(b) 时刻 $t$ 正在传回发送方的在途 ACK 取值范围是 $[k-5,\; k-1]$（即 $\{k-5,\; k-4,\; k-3,\; k-2,\; k-1\} \pmod{1024}$，共 5 个可能值）}：
  \begin{itemize}
  \item \textbf{为什么 $[k-4, k-1]$ 可以在途}：接收方刚刚收妥 $k-4, k-3, k-2, k-1$ 时依次生成的这 4 个确认报文，在时刻 $t$ 完全可能仍在回程链路上飞行尚未到达发送方。
  \item \textbf{为什么 $\mathrm{ACK}(k-5)$ 也完全可以在无乱序信道上处于在途状态（极易漏算的第 5 个值）}：虽然让发送方滑窗到 $[k-4, k-1]$ 的\emph{第一个} $\mathrm{ACK}(k-5)$ 早已抵达发送方，但接收方在等待 $k-4$ 期间若收到失序高号分组，还会\textbf{再次生成新的重复确认 $\mathrm{ACK}(k-5)$}。构造固定单程时延 $D$ 的无乱序反例：设发送方在收到首个 $\mathrm{ACK}(k-5)$ 前已发出 $\mathrm{pkt}(k-4)$（途中丢失）和紧随其后的 $\mathrm{pkt}(k-2)$；稍后发送方对 $\mathrm{pkt}(k-4)$ 超时，紧跟在 $\mathrm{pkt}(k-2)$ 身后重传 $\mathrm{pkt}(k-4)..\mathrm{pkt}(k-1)$。当旧的 $\mathrm{pkt}(k-2)$ 先抵达仍在等待 $k-4$ 的接收方时，接收方丢弃它并向回程链路发出一个\textbf{重复确认 $\mathrm{ACK}(k-5)$}；紧接着重传的 $\mathrm{pkt}(k-4)..\mathrm{pkt}(k-1)$ 鱼贯抵达接收方，使接收方在时刻 $t$ 的期望序号顺利推进至 $k$——此刻回程链路上最前面飞行的恰恰就是刚才那枚尚未到达发送方的重复 $\mathrm{ACK}(k-5)$，身后跟着 $\mathrm{ACK}(k-4)..\mathrm{ACK}(k-1)$。
  \item \textbf{为什么严格小于 $k-5$ 的 ACK 不可能在途}：一旦发送方早已收到了首个 $\mathrm{ACK}(k-5)$，在信道不乱序的前提下，所有在此之前发出的更旧确认（$\le k-6$）必定早已在首个 $\mathrm{ACK}(k-5)$ 之前到达（或丢失），且接收方自收妥 $k-5$ 后再也不会生成小于 $k-5$ 的确认号；同时接收方尚未收到 $k$，不可能发出 $\ge k$ 的确认号。故在途 ACK 取值集合严格为 $\{k-5, k-4, k-3, k-2, k-1\} \pmod{1024}$。
  \end{itemize}
\end{enumerate}

\textbf{English answer:} (a) With $N=4$, the sender's window has size 4 and can start anywhere in $[k-4, k]$, giving the 5 possible windows $[k-4,k-1], [k-3,k], [k-2,k+1], [k-1,k+2], [k,k+3] \pmod{1024}$. (b) Propagating ACK values can take any of the \textbf{5 values} in $[k-N-1, k-1] = \{k-5, k-4, k-3, k-2, k-1\} \pmod{1024}$: ACKs for $k-4..k-1$ may still be in flight, and a duplicate $\mathrm{ACK}(k-5)$ triggered just before packet $k-4$ arrived may also still be propagating back; no ACK $< k-5$ can remain in transit because the sender already received an $\mathrm{ACK}(k-5)$ before sending $k-1$ (and the channel does not reorder), and no ACK $\ge k$ has been sent yet.


\StudyHeading{1}{Ch.4 \& 5 P4. 传统数据报转发表与“按源主机分流”限制}
\StudyTarget{CS5222-Assignment02__p4}

\textbf{Original Question:} Consider the datagram network shown below (left panel of Figure): Router A connects to H1 on interface 1, H2 on interface 2, Router B on interface 3, and Router C on interface 4; both B and C connect to Router D, which connects to H3.
\begin{itemize}
\item \textbf{(a)} Show the forwarding table in Router A, such that all traffic destined to host H3 is forwarded through interface 3.
\item \textbf{(b)} Can you write down a forwarding table in Router A, such that all traffic from H1 destined to host H3 is forwarded through interface 3, while all traffic from H2 destined to host H3 is forwarded through interface 4? (Hint: this is a trick question.)
\end{itemize}

\StudyImage{CS5222/course-notes/assets/ass2-topologies.png}
\StudyCaption{原创重绘 Assignment 2 两道网络拓扑：左图为 P4 传统数据报网络（含各路由器接口编号）；右图为 P19–P20 的三交换机 SDN OpenFlow 网络（完整标出主机 IP 与交换机端口 1–4）。}

\textbf{答案与解析：}
\begin{itemize}
\item \textbf{(a) 路由器 A 的转发表}：\texttt{Destination Address = H3  -->  Link Interface = 3}。
\item \textbf{(b) 不能（No）}：传统 IP 数据报网络的转发表\textbf{仅根据分组的目的 IP 地址（destination address only）}做最长前缀匹配查表，不检查源 IP 地址（H1 还是 H2）或入接口编号；既然两股流量的目的地址都是 H3，传统路由器 A 无法将它们导向不同出接口（必须使用下面 P19–P20 的 SDN OpenFlow 广义流表才能实现）。
\end{itemize}

\textbf{English answer:} (a) Forwarding table at A: \texttt{Destination = H3 | Interface = 3}. (b) \textbf{No}, because forwarding rules in a traditional destination-based datagram network match \emph{only} on the destination IP address and cannot distinguish packets with the same destination (H3) by their source host (H1 vs. H2).


\StudyHeading{1}{Ch.4 P19. OpenFlow 交换机 \texttt{s2} 流表编写（双向互通与本地交付）}
\StudyTarget{CS5222-Assignment02__p19}

\textbf{Original Question:} Consider the SDN OpenFlow network shown in the right panel above. Specify the flow table entries in \textbf{\texttt{s2}} that implement the following desired forwarding behavior for datagrams arriving at \texttt{s2}:
\begin{itemize}
\item any datagrams arriving on input port 1 from hosts \texttt{h5} or \texttt{h6} that are destined to hosts \texttt{h1} or \texttt{h2} should be forwarded over output port 2;
\item any datagrams arriving on input port 2 from hosts \texttt{h1} or \texttt{h2} that are destined to hosts \texttt{h5} or \texttt{h6} should be forwarded over output port 1;
\item any arriving datagrams on input ports 1 or 2 and destined to hosts \texttt{h3} or \texttt{h4} should be delivered to the host specified;
\item hosts \texttt{h3} and \texttt{h4} should be able to send datagrams to each other.
\end{itemize}

\textbf{核对 \texttt{s2} 的四个物理端口映射（见上方右图）：}Port 1 连 \texttt{s3}（\texttt{10.3.0.*}：含 \texttt{h5=10.3.0.5, h6=10.3.0.6}），Port 2 连 \texttt{s1}（\texttt{10.1.0.*}：含 \texttt{h1=10.1.0.1, h2=10.1.0.2}），Port 3 连本地主机 \texttt{h3 (10.2.0.3)}，Port 4 连本地主机 \texttt{h4 (10.2.0.4)}。按课件与官方参考答案的标准单端口匹配格式，\texttt{s2} 的完整 8 条流表项如下：

\begingroup\small\setlength{\tabcolsep}{3pt}\renewcommand{\arraystretch}{1.20}
\begin{longtable}{@{}>{\raggedright\arraybackslash}p{\dimexpr0.08\linewidth-1.33333\tabcolsep\relax}>{\raggedright\arraybackslash}p{\dimexpr0.66\linewidth-1.33333\tabcolsep\relax}>{\raggedright\arraybackslash}p{\dimexpr0.26\linewidth-1.33333\tabcolsep\relax}@{}}
\toprule
行号 & Match Fields（未列出字段均为通配符 \texttt{*}） & Action \\
\midrule\endfirsthead
行号 & Match Fields（未列出字段均为通配符 \texttt{*}） & Action \\
\midrule\endhead
\bottomrule\endfoot
1 & \texttt{Ingress Port = 1; IP Src = 10.3.0.*; IP Dst = 10.1.0.*} & \texttt{Forward(2)} \\
2 & \texttt{Ingress Port = 2; IP Src = 10.1.0.*; IP Dst = 10.3.0.*} & \texttt{Forward(1)} \\
3 & \texttt{Ingress Port = 1; IP Dst = 10.2.0.3} & \texttt{Forward(3)} \\
4 & \texttt{Ingress Port = 2; IP Dst = 10.2.0.3} & \texttt{Forward(3)} \\
5 & \texttt{Ingress Port = 1; IP Dst = 10.2.0.4} & \texttt{Forward(4)} \\
6 & \texttt{Ingress Port = 2; IP Dst = 10.2.0.4} & \texttt{Forward(4)} \\
7 & \texttt{Ingress Port = 3; IP Src = 10.2.0.3; IP Dst = 10.2.0.4}（或简写 \texttt{Ingress Port = 3}） & \texttt{Forward(4)} \\
8 & \texttt{Ingress Port = 4; IP Src = 10.2.0.4; IP Dst = 10.2.0.3}（或简写 \texttt{Ingress Port = 4}） & \texttt{Forward(3)} \\
\end{longtable}\endgroup
注：因为属于 \texttt{h5/h6} 的子网仅占 \texttt{10.3.0.5, 10.3.0.6}、\texttt{h1/h2} 仅占 \texttt{10.1.0.1, 10.1.0.2}，表格中的子网通配符 \texttt{10.3.0.*} 与 \texttt{10.1.0.*} 在该实验网内与逐一展开主机 IP 等价；若题目禁止子网通配符，将第 1、2 行拆成显式主机对即可。


\StudyHeading{1}{Ch.4 P20. OpenFlow 交换机 \texttt{s2} 顺时针与逆时针按源分流}
\StudyTarget{CS5222-Assignment02__p20}

\textbf{Original Question:} Consider again the SDN OpenFlow network above. Specify the flow table entries in \textbf{\texttt{s2}} for datagrams arriving from hosts \texttt{h3} or \texttt{h4} at \texttt{s2} such that:
\begin{itemize}
\item any datagrams arriving from host \texttt{h3} and destined for \texttt{h1, h2, h5}, or \texttt{h6} should be forwarded in a \textbf{clockwise} direction in the network;
\item any datagrams arriving from host \texttt{h4} and destined for \texttt{h1, h2, h5}, or \texttt{h6} should be forwarded in a \textbf{counter-clockwise} direction in the network.
\end{itemize}

\textbf{几何方向核对与精确四行流表：}观察上方右图的三角形环（左上 \texttt{s3}、右上 \texttt{s2}、正下 \texttt{s1}），站在右上角的 \texttt{s2} 往外发包：\textbf{顺时针方向}是向右下前往 \texttt{s1}（对应 \texttt{s2} 的 \textbf{出端口 2}）；\textbf{逆时针方向}是向左前往 \texttt{s3}（对应 \texttt{s2} 的 \textbf{出端口 1}）。为避免把 \texttt{h3 <-> h4} 本地互发流量误送入环网，目的前缀明确匹配 \texttt{10.1.0.*} 与 \texttt{10.3.0.*}：

\begingroup\small\setlength{\tabcolsep}{3pt}\renewcommand{\arraystretch}{1.20}
\begin{longtable}{@{}>{\raggedright\arraybackslash}p{\dimexpr0.08\linewidth-1.33333\tabcolsep\relax}>{\raggedright\arraybackslash}p{\dimexpr0.64\linewidth-1.33333\tabcolsep\relax}>{\raggedright\arraybackslash}p{\dimexpr0.28\linewidth-1.33333\tabcolsep\relax}@{}}
\toprule
行号 & Match Fields in \texttt{s2} & Action \\
\midrule\endfirsthead
行号 & Match Fields in \texttt{s2} & Action \\
\midrule\endhead
\bottomrule\endfoot
1 & \texttt{Ingress Port = 3; IP Src = 10.2.0.3; IP Dst = 10.1.0.*} & \textbf{\texttt{Forward(2)}}（顺时针往 \texttt{s1}） \\
2 & \texttt{Ingress Port = 3; IP Src = 10.2.0.3; IP Dst = 10.3.0.*} & \textbf{\texttt{Forward(2)}}（顺时针往 \texttt{s1}） \\
3 & \texttt{Ingress Port = 4; IP Src = 10.2.0.4; IP Dst = 10.1.0.*} & \textbf{\texttt{Forward(1)}}（逆时针往 \texttt{s3}） \\
4 & \texttt{Ingress Port = 4; IP Src = 10.2.0.4; IP Dst = 10.3.0.*} & \textbf{\texttt{Forward(1)}}（逆时针往 \texttt{s3}） \\
\end{longtable}\endgroup

\textbf{独立变式 / Check：}若在 P19 基础上增加一条安全策略——禁止 \texttt{h5/h6} 访问 \texttt{h3} 的 SSH 服务（TCP 目的端口 22），其余流量照常放行，应在 \texttt{s2} 增加怎样一条最高优先级 OpenFlow 表项？ / Write a high-priority OpenFlow rule at \texttt{s2} to block TCP port-22 traffic from \texttt{h5/h6} to \texttt{h3}.

\begin{StudySolution}{答案 / Answer}
在 \texttt{s2} 顶部加入高优先级防火墙匹配项：\texttt{Ingress Port = 1; IP Src = 10.3.0.*; IP Dst = 10.2.0.3; IP Prot = 6 (TCP); TCP Dst Port = 22  -->  Action: Drop}（无转发动作即丢弃）。 / High-priority entry: \texttt{Ingress Port=1, IP Src=10.3.0.*, IP Dst=10.2.0.3, IP Prot=6, TCP dport=22 -> Drop}.
\end{StudySolution}

来源：\texttt{CS5222 Assignment 2.pdf}（3 页）与课程标准参考解答对照。
